\appendix
\section{Notation and interpretation glossary}
\label{app:glossary}

\begin{table}[h]
\centering
\begin{tabular}{@{}p{0.27\linewidth}p{0.69\linewidth}@{}}
\toprule
Symbol or term & Meaning\\
\midrule
$a_i$ & Candidate $i$: a complete 30-step physical action sequence.\\
Endpoint $s_i$ & State after executing all actions in candidate $i$.\\
$z_i$, $\hat z_i$, $z_g$ & Encoded realized endpoint, predicted terminal latent, encoded goal.\\
$C_{\mathrm{pred}}$, $C_{\mathrm{real}}$ & The same latent goal cost evaluated on predicted or encoded realized endpoints.\\
$a_{\mathrm{pred}}$ & Candidate selected by predicted latent cost.\\
$a_{\mathrm{real}}$ & Candidate selected by realized-endpoint latent cost; retrospective, not a physical optimum.\\
$a_{\mathrm{state}}$ & Candidate minimizing realized native terminal distance; retrospective and non-deployable.\\
$D$, $U_{\mathrm{dist}}$ & Native distance and its negative, the higher-is-better distance utility.\\
$R_{\mathrm{task}}$ & Cumulative official Push-T coverage reward or Wall binary completion.\\
Record & One model, initial-state/goal instance, and proposal seed, with 300 candidates.\\
Paired cell & One Push-T trajectory offset and seed with valid records for both models.\\
Cluster & A scheduled Push-T trajectory or a complete Wall case; the resampling unit.\\
Initial candidates & Samples before any CEM distribution update; historically called an ``iteration-0 menu'' in code.\\
\bottomrule
\end{tabular}
\end{table}

\section{Original inference, unchanged}
\label{app:original}

Table~\ref{tab:original} retains the original pooled point estimates, percentile
intervals, and Holm-adjusted sign-flip values. Push-T uses all 20 scheduled
trajectory contrasts, 10,000 trajectory-bootstrap draws, and enumeration of
all $2^{20}$ sign assignments for its two-sided paired test. Wall uses 29
complete-case contrasts, 10,000 case-bootstrap draws (seed 20260911), and
100,000 random sign assignments with the $+1$ correction (seed 20260912).
Holm adjustment covers two outcomes separately for endpoint correction and
predicted-minus-uniform selection within each environment. Push-T's bounded
candidate sensitivity has its own originally declared two-outcome family.
The exact test requires sign symmetry of cluster contrasts under the null.
No new state-selector comparison has been added to those original families.

\begin{table}[h]
\caption{Unchanged original inference. Intervals are marginal 95\% cluster
bootstrap intervals; $p_H$ is the original within-family Holm-adjusted value.}
\label{tab:original}
\centering\small
\begingroup
\setlength{\tabcolsep}{3pt}
\begin{tabular}{lllrr}
\toprule
Environment & Contrast & Outcome & Effect [95\% CI] & $p_H$\\
\midrule
Push-T & Realized $-$ predicted & Distance utility & $-5.299\ [-9.635,-0.706]$ & .075539\\
 & & Reward & $+0.04183\ [-0.17569,0.24024]$ & .701128\\
 & Predicted $-$ uniform & Distance utility & $+48.702\ [40.473,56.434]$ & $7.6294\!\times\!10^{-6}$\\
 & & Reward & $+0.76099\ [0.42033,1.11714]$ & .00040245\\
Wall & Realized $-$ predicted & Distance utility & $-0.25365\ [-0.90933,0.36436]$ & .923091\\
 & & Completion & $+0.02874\ [0,0.08046]$ & .923091\\
 & Predicted $-$ uniform & Distance utility & $+12.36271\ [10.90015,13.75461]$ & $1.99998\!\times\!10^{-5}$\\
 & & Completion & $+0.01705\ [0,0.04544]$ & .501395\\
\bottomrule
\end{tabular}
\endgroup
\end{table}

The new diagnostic uses seeds 20261002 (Push-T) and 20261003 (Wall), with
10,000 whole-cluster draws shared by its three contrasts and both outcomes
within each environment. These intervals quantify the new post-hoc comparisons
and do not replace the original test decisions. Recomputing the original
results is a required preflight check before generating the new results.

\section{Model-specific descriptive comparisons}
\label{app:models}

Table~\ref{tab:models} uses the same retained cells, cluster definitions,
denominators, and common bootstrap draws as the pooled post-hoc analysis.
Each model estimate compares a latent selector with the same physical state
selector on matched sets. Repeated state-selected outcomes do not provide a
second physical experiment. All entries are descriptive; differences between
their confidence intervals are not model-interaction tests.

\begin{table}[h]
\caption{Post-hoc model-specific effects with 95\% cluster-bootstrap
intervals. ``Real'', ``pred'', and ``state'' refer to the three defined selectors.}
\label{tab:models}
\centering\small
\begin{tabular}{lllrr}
\toprule
Task & Model & Contrast & $\Delta U_{\mathrm{dist}}$ [95\% CI] & $\Delta R_{\mathrm{task}}$ [95\% CI]\\
\midrule
Push-T & JEPA-WM & Real $-$ pred & $-9.330\ [-14.394,-4.557]$ & $+.081\ [-.130,.280]$\\
 & & State $-$ real & $+67.623\ [59.887,76.057]$ & $-.861\ [-1.271,-.453]$\\
 & & State $-$ pred & $+58.293\ [50.879,65.943]$ & $-.780\ [-1.169,-.407]$\\
 & DINO-WM & Real $-$ pred & $-1.268\ [-7.898,6.272]$ & $+.003\ [-.241,.237]$\\
 & & State $-$ real & $+72.005\ [63.922,80.083]$ & $-.815\ [-1.179,-.462]$\\
 & & State $-$ pred & $+70.737\ [64.052,77.512]$ & $-.812\ [-1.229,-.420]$\\
\addlinespace
Wall & JEPA-WM & Real $-$ pred & $-.655\ [-1.649,.215]$ & $+.0345\ [0,.0920]$\\
 & & State $-$ real & $+4.910\ [3.327,6.785]$ & $0\ [0,0]$\\
 & & State $-$ pred & $+4.255\ [3.001,5.839]$ & $+.0345\ [0,.0920]$\\
 & DINO-WM & Real $-$ pred & $+.148\ [-.724,.923]$ & $+.0230\ [0,.0690]$\\
 & & State $-$ real & $+6.376\ [4.222,8.825]$ & $0\ [0,0]$\\
 & & State $-$ pred & $+6.524\ [4.551,8.761]$ & $+.0230\ [0,.0690]$\\
\bottomrule
\end{tabular}
\end{table}

\section{Sampling, missingness, and the recovery amendment}
\label{app:protocol}

\paragraph{Push-T selection and support.}
The manifest selects validation trajectory indices 0--20 except 7, for 20
trajectories, without consulting audit outcomes. For a trajectory of length
$L$, its five offsets are $\lfloor k(L-31)/4\rfloor$, $k=0,\ldots,4$.
Only recorded sequence lengths select the offsets. Expert block-state motion
labels a secondary phase proxy but does not select or exclude records.
Each offset is paired with proposal seeds 1, 2, and 3 and both models.
The resulting 100 distinct initial-state/goal segments produce 300 offset--seed
cells and 600 scheduled model records. The manifest's historical
\texttt{segment\_count=300} includes proposal seeds.

For an outcome $Y$ (either $U_{\mathrm{dist}}$ or $R_{\mathrm{task}}$), let
$\delta_{tcm}$ be a specified selector contrast for trajectory $t$, cell $c$,
and model $m$. The cell indicator $I_{tc}$ equals one only when both model
records are valid under the original rules. The pooled trajectory contrast is
\begin{equation}
 \bar\delta_t=\frac{1}{2}\sum_{m=1}^{2}
          \frac{1}{15}\sum_{c=1}^{15} I_{tc}\delta_{tcm},
 \qquad \widehat\Delta=\frac{1}{20}\sum_{t=1}^{20}\bar\delta_t.
 \label{eq:cluster}
\end{equation}
Missing-cell terms are defined as zero \emph{contrast contributions}; the
formula does not assign a physical outcome to a missing action. The two-model
mask and fixed 15-cell denominator are preserved for every new contrast and
model split. Observed-cell means would answer a different question.
The original support requirement is at least 540 matched model records and
at least 18 trajectories with 12 or more matched cells. The retained study
has 562 matched records and 19 such trajectories, so it meets that requirement;
all 20 scheduled trajectories still enter Eq.~\ref{eq:cluster}.
Support gates require coverage of the scheduled experiment, not a positive
physical effect or a guarantee of adequate power. The original broad-recovery
rule additionally required passed integrity and repeat-stability checks;
positive pooled means and lower 95\% bounds, with Holm-adjusted $p<.05$, for
both outcomes under raw and bounded candidates; positive model-specific lower
bounds for both outcomes in both candidate sets; and no conflicting pooled
signs. Failure of any condition blocked the broad claim. Sampling, outcomes,
aggregation, tests, and this rule were recorded in the original locked protocol;
recovery was a later amendment, and the auxiliary and new state-selector
analyses are separately labeled post-hoc.

\paragraph{Infrastructure recovery.}
The Push-T protocol was recorded as locked on August 30, 2026 and prohibited
retries. Its original 600-record attempt ledger ended with 420 completed and
180 failed records. A fatal CUDA error had contaminated the persistent worker,
causing later untouched records to receive the same infrastructure error.
The recovery amendment mechanically selected 169 failed records with that
exact error message for at most one attempt each, in isolated processes.
The remaining failure classes were not eligible. The amendment states that
eligibility was fixed before inspecting scientific outcomes, preserving
candidate specifications and analysis rules, but this remains a deviation
from the original no-retry rule. Original failures remain in the ledger;
recovery neither deletes them nor makes them original completions.
The no-retry rule restricted opportunities to change retained evidence through
selective repeated attempts; the amendment does not retroactively satisfy it.
The recovery-augmented analysis has
577 valid and 23 invalid records, with 15 valid records unmatched across models.
Removing every recovery-sourced artifact leaves 408 matched records and only
14 trajectories meeting the original eligibility threshold: the original-only
subset fails the original support requirement. Its distance and reward effects
are $-5.39$ ($[-9.32,-1.40]$) and $-0.019$ ($[-0.212,0.160]$), respectively.
These are exploratory sensitivity estimates, not a restored confirmatory
analysis or evidence that missingness is ignorable. The 15 observed trajectory
groups used by a separate diagnostic are not 15 eligible primary clusters.

Earlier protocol versions are not scientific inputs to this analysis: v1.2
failed its prescribed RGB/serialization integrity check, and v1.3 failed
before usable measurements because a replay-encoding helper lacked its
required model argument. The v1.4 source and evidence are separately bound;
no earlier-version result is reused. The retained protocol, amendment, and
ledgers make this history inspectable but do not constitute an independently
registered timestamp for the declared outcome-blind sequence.

\paragraph{Wall selection and missingness.}
Case seeds 2026091100--2026091129 generate the 30 initial-state/goal cases.
Engineering seeds 2026091001 and 2026091002 are excluded from the study.
The wall location is 32 and door location 30 in the underlying 65-pixel
environment. Each case contains seeds 1--3 and both models. One JEPA-WM
record, case \texttt{wall\_028}, seed 1, failed scientific validation and was
not retried. Its entire case is excluded from the complete-case physical
analysis. The remaining five valid records contribute only to all-valid
descriptions. Case contrasts average the six model--seed records equally;
the 29 case means are the inferential units. This restriction does not rule
out difficulty-dependent missingness.
The retained failure ledger labels this failure \texttt{scientific\_or\_child}
but supplies no error message, so its specific failed criterion is not
established by that ledger.

\paragraph{Integrity and ties.}
Original Push-T validity required finite production arrays, decomposition
residuals at most $10^{-8}$, and score reproduction within $2\times10^{-5}$
with zero relative tolerance. Re-encoding serialized observations at chunk
sizes eight and four had to preserve scores within that tolerance and the
latent minimum tie sets exactly. Original repeat and provenance checks also
applied. Execution completion alone therefore did not imply scientific
validity; a valid record still needed its partner for paired analysis.
The new preflight checks cover all 780 unique scheduled records: 756 valid
and 24 invalid across both environments. They verify candidate counts,
selectors, stored physical arrays, cross-model agreement, and the original
headline analysis before new computation. Push-T's original latent selectors
require a unique minimum under absolute tolerance $10^{-10}$ and relative
tolerance $10^{-8}$; otherwise that record is invalid. Wall's latent selectors
use the first exact minimum. The new state selector also uses the first exact
minimum. No exact minimum ties or duplicate physical sequences occur in the
retained valid candidate sets; synthetic tests cover deterministic tie
handling. Cross-model identity is established by arrays, not assumed from
shared random-seed labels.

\section{Model and environment configuration}
\label{app:config}

\paragraph{Visual encoding.}
Both models in both tasks load \texttt{dinov2\_vits14} from
\texttt{dinov2\_vits14\_pretrain.pth}: a ViT-S/14 with 384-dimensional patch
features. Raw RGB is divided by 255 and resized to $224\times224$ using
bilinear interpolation with \texttt{align\_corners=False}. The audit loaders
fix spatial scale and aspect-ratio augmentation ranges to $(1,1)$ and disable
flips, motion augmentation, automatic augmentation, and random erasing.
DINO-WM then uses per-channel mean and standard deviation $(0.5,0.5,0.5)$.
JEPA-WM uses mean $(0.485,0.456,0.406)$ and standard deviation
$(0.229,0.224,0.225)$. Sharing backbone weights therefore does not mean that
the models receive identically normalized inputs. Each model's preprocessing
is fixed between its predicted and realized costs.

\paragraph{Proprioception and actions.}
Push-T proprioception contains the agent's two-dimensional position and
velocity (four coordinates), not block pose. Wall proprioception is the
two-dimensional agent position. After the released coordinate normalization,
the frozen proprioceptive encoder is a kernel-one, stride-one convolution,
equivalent to a per-time-step affine map. Its output dimension is 20 for
DINO-WM on Push-T and 10 on Wall, versus 16 for JEPA-WM on both tasks.
Table~\ref{tab:normalization} gives the environment normalization constants;
the operational meaning of action normalization is
$a_{\mathrm{physical}}=\sigma_a\odot a_{\mathrm{normalized}}+\mu_a$.
Consequently, candidate zero is a normalized-mean sequence, not necessarily
a physical no-op. Gaussian proposals are unbounded in normalized coordinates;
the bounded sensitivity is defined separately, not silently substituted.

\begin{table}[h]
\caption{Released environment normalization constants, rounded here for
readability. Full-precision values remain in the bound dataset registry.}
\label{tab:normalization}
\centering\small
\begin{tabular}{@{}p{0.10\linewidth}p{0.32\linewidth}p{0.50\linewidth}@{}}
\toprule
Task & Quantity & Values\\
\midrule
Push-T & Action mean & $(-0.0087,\ 0.0068)$\\
 & Action standard deviation & $(0.2019,\ 0.2002)$\\
 & Proprioception mean & $(236.61549,\ 264.56741,\ -2.93032,\ 2.54308)$\\
 & Proprioception standard deviation & $(101.12020,\ 87.01120,\ 74.84556,\ 74.14009)$\\
Wall & Action mean & $(-0.00556462,\ 0.01319782)$\\
 & Action standard deviation & $(0.75233287,\ 0.74668980)$\\
 & Proprioception mean & $(31.89451,\ 31.81402)$\\
 & Proprioception standard deviation & $(17.66352,\ 17.27517)$\\
\bottomrule
\end{tabular}
\end{table}

The model horizon is six, action dimension ten, physical action dimension two,
and frame grouping five, giving 30 executed actions with
\texttt{action\_skip=1}. Initial image and proprioceptive history lengths are
one; rollout context is two. The objective is terminal visual MSE plus
$0.1$ proprioceptive MSE, with \texttt{sum\_all\_diffs=False}.
Push-T samples with a PyTorch generator on the initial latent's device;
Wall uses a CPU generator with shape $[6,300,10]$ before permutation.
Both use standard-normal proposals, but their random streams should not be
assumed identical across environments.

\paragraph{Dataset and goals.}
Push-T uses \path{data/pusht_noise/val} in the retained official source tree.
The simulator starts from \texttt{states[offset]} with the recorded shape,
and replays the 30 recorded relative actions using environment seed 1 to
generate the goal. The goal is not merely a frame extracted from the dataset
video. Wall cases are generated from the fixed random-state seeds described
above rather than validation-trajectory offsets. The capture routine
\texttt{run\_case} in \path{experiments/wall_decision_audit.py} calls the released
\texttt{WallEnvWrapper.sample\_random\_init\_goal\_states}, which delegates to
\texttt{DotWall.generate\_random\_state} in
\path{evals/simu_env_planning/envs/wall_env/envs/wall.py} under the bound
Wall configuration. It samples positions uniformly on opposite sides of the
wall within padded boundaries, reversing travel direction with probability
one half. These routines, not the seed alone, define case generation.

\paragraph{Source and checkpoint bindings.}
The official JEPA-WM source commit is
\nolinkurl{13cf1d9c7e476f53c17714d2e0f1dc239a883ce0};
the DINOv2 source commit is
\nolinkurl{b194f00db6136677fc8a4cc2ef2168f7699dfba2}.
The recorded Wall model-repository revision is
\nolinkurl{9b9c41ef249466630dbf1a20e78391865d07b3b9}.
The checkpoint filenames and SHA-256 digests are:
\begin{itemize}
\item \texttt{dino\_wm\_pusht.pth.tar}:\\
\texttt{8ec9cb05f22812d7f12e3c216b0637f4}\allowbreak\texttt{1641055c0653e503e2746edb981b550f}.
\item \texttt{jepa\_wm\_pusht.pth.tar}:\\
\texttt{9beca3eafe0739c3b3adb5d734fa435c}\allowbreak\texttt{cbda0fea8a65d53d4cccec176aaaa0eb}.
\item \texttt{dino\_wm\_wall.pth.tar}:\\
\texttt{ff170be5aec9249768be4a220d600b8f}\allowbreak\texttt{00a8589b2a78982ecf9273809f2767df}.
\item \texttt{jepa\_wm\_wall.pth.tar}:\\
\texttt{8efb0623cfba1cb3ca210de26f7579c8}\allowbreak\texttt{3dd24936635f11989c515afcb23bea1e}.
\end{itemize}
The DINOv2 weight-file digest is
\texttt{b938bf1bc15cd2ec0feacfe3a1bb553f}\allowbreak\texttt{e8ea9ca46a7e1d8d00217f29aef60cd9};
the recorded loaded state-dictionary digest is
\texttt{5bb315664719c85935c4b67a0f8faef9}\allowbreak\texttt{d01b7dfb57ffbecef1c062f7b4e4aa35}.
The Push-T dataset archive digest is
\texttt{442f5dee246edf670964ed7bdecd2486}\allowbreak\texttt{83cd6d00580fa0e4d458abb53f92da08}.

These facts are bound by
\path{experiments/official_pusht_v14_manifest.json}
(\texttt{selection}, \texttt{protocol}, \texttt{source\_bundle});
\path{output/source/physical_ai_supplement/official_pusht_v14_certificate_issuance.json}
(model, encoder, and data bindings);
\path{experiments/wall_audit_20260910/protocol.json};
\path{experiments/wall_audit_20260910/execution_bindings.json}; and
\path{output/wall_recompute_20260910_v3/metadata/provenance.json}.
The official planning configurations lie under
\path{configs/evals/simu_env_planning/}, with task directories \texttt{pt}
and \texttt{wall}, model directories \texttt{dino-wm} and \texttt{jepa-wm},
and names ending
\path{L2_cem_sourcedset_H6_nas6_ctxt2_r224_alpha0.1_ep96_decode.yaml}
for Push-T or the analogous \texttt{sourcerandstate} name for Wall.
Normalization constants are in
\path{app/plan_common/datasets/__init__.py} of the bound source tree.
Checkpoint training epochs, training seeds, and complete training-data
provenance are not established by these retained audit bindings; we do not
infer them from filenames.

\section{Earlier sensitivities and exact score accounting}
\label{app:sensitivity}

These analyses predate the state-selector extension and retain their original
status and numerical summaries. They are not new experiments for this revision.
The bounded-candidate and original-completions-only sensitivities concern
endpoint correction, not the new state-selector contrasts. We have not
reported an observed-cell-denominator analysis or an original-completions-only
state-selector analysis. Their absence leaves the corresponding robustness
questions unresolved; the existing sensitivities cannot establish that
missingness is ignorable for any contrast.

\paragraph{Bounded candidates.}
Restrict each set to sequences with maximum absolute normalized action
coordinate at most three. The original support requirement is at least 200
candidates and valid untied latent selectors. The realized-minus-predicted
effects are $-5.80$ in distance utility (95\% CI $[-11.00,-0.21]$,
Holm-adjusted $p=.111$) and $+0.115$ in reward ($[-0.093,0.314]$, $p=.296$).
The raw analysis remains primary; this sensitivity does not alter the new
state-selector definition or its candidate set.

\paragraph{Physical components.}
The component analysis was post-hoc, specified before inspecting component
outcomes. It uses the same 562 matched records, 281 cells, 20 trajectories,
and fixed denominators. Position and velocity utilities are negative
Euclidean distances to their goal components; angle utility is negative
circular absolute error, unlike the raw angle coordinate in native distance.
Terminal coverage reward is the last official reward, not the raw overlap
fraction reported separately by the environment. Holm adjustment covers these
five declared component tests.

\begin{table}[h]
\caption{Retained post-hoc Push-T component effects, realized minus predicted.
Intervals are original marginal 95\% trajectory-bootstrap intervals.}
\label{tab:appendix_components}
\centering
\begin{tabular}{lrrr}
\toprule
Component utility & Effect & 95\% CI & $p_H$\\
\midrule
Agent position & $-1.817$ & $[-4.564,.971]$ & .232\\
Block position & $+1.267$ & $[.480,2.176]$ & .0203\\
Circular block angle & $+.0325$ & $[.0197,.0471]$ & .000210\\
Agent velocity & $-5.809$ & $[-10.523,-1.187]$ & .0957\\
Terminal coverage & $+.00901$ & $[.00035,.01696]$ & .110\\
\bottomrule
\end{tabular}
\end{table}

For exact native squared-distance accounting, the mean predicted-selection
minus realized-selection contributions are $-477.2$ from agent position,
$+79.2$ from block position, $+0.226$ from raw angle, and $-1497.0$ from agent
velocity. Positive contributions mean realized selection reduces that squared
error. The maximum rowwise identity residual is $3.64\times10^{-11}$.
These mixed-unit squared contributions are not additive contributions to mean
unsquared distance utility, and they change under coordinate rescaling.

\paragraph{Stored-latent interpolation.}
For each candidate, the earlier post-hoc sensitivity scores
$(1-\lambda)\hat z_i+\lambda z_i$, for $\lambda=0,.1,\ldots,1$, with all
actions, goals, cost function, and physical outcomes otherwise unchanged. At
$\lambda=.5$, 38.4\% of matched-record choices change and the distance effect
is $-2.24$; at $\lambda=1$, these values are 58.9\% and $-5.30$.
All ten nonzero-$\lambda$ pooled distance point estimates are negative.
These aggregate points are not independent replicates, interpolated latents
need not represent realizable states, and no trainable correction is evaluated.

\paragraph{Score accounting is an identity, not a causal explanation.}
For either modality, write its prediction error as $e_i=\hat z_i-z_i$, with
dimension $d$. Then
\begin{equation}
 \frac{\|\hat z_i-z_g\|_2^2-\|z_i-z_g\|_2^2}{d}
 =\frac{2(z_i-z_g)^\top e_i}{d}+\frac{\|e_i\|_2^2}{d}.
 \label{eq:identity}
\end{equation}
Summing the two modalities with weights one and $0.1$ reconstructs
$C_{\mathrm{pred}}(i)-C_{\mathrm{real}}(i)$ exactly; subtracting two candidates'
identities reconstructs their margin change. The maximum residual in retained
Push-T records is $6.66\times10^{-16}$. Directional terms dominate magnitude
terms in 266 of 337 selector switches (78.9\%): the absolute directional
margin between the predicted and realized winners exceeds the absolute
magnitude margin after the visual and $0.1$ proprioceptive contributions are
combined. This is a record-level fraction, not cluster-level inference.
This explains algebraic score
movement, not its physical desirability or a causal failure mechanism.
The trajectory-weighted mean candidate-set Spearman correlation between the
two costs is $.858$ (95\% CI $[.839,.875]$); the mean per-set median absolute
cost difference is $.0505$, with mean ratio $.339$ to predicted-cost IQR.

\paragraph{Ranking summaries.}
Spearman correlates negative cost with higher-is-better outcome, using average
ranks for exact ties. Pairwise accuracy counts whether lower cost favors the
higher outcome, excluding pairs tied on either quantity. Per-record summaries
use the original paired-cell, equal-model, trajectory aggregation.
On bounded candidate support, the earlier descriptive distance-utility
Spearman correlations are .367 for predicted and .345 for realized scoring;
pairwise ranking accuracies are .599 and .591. Their normalized regrets are
.250 and .274. For cumulative reward, the corresponding pairs are .121/.124,
.469/.471, and .372/.359. For a higher-is-better outcome $Y$, normalized
regret of selection $a$ is
$(\max_iY(a_i)-Y(a))/(\max_iY(a_i)-\min_iY(a_i))$ when its denominator is
positive and the selector is defined. These summaries are descriptive and
do not establish ranking superiority.

\section{Secondary selector-instability diagnostic}
\label{app:neartie}

The earlier near-tie diagnostic targets whether $a_{\mathrm{pred}}$ and
$a_{\mathrm{real}}$ differ, not whether either fails physically. After
deduplicating physical actions and retaining the lowest predicted cost per
action, let $c_{(1)}$ and $c_{(2)}$ be the two lowest predicted costs.
The ambiguity score is
\begin{equation}
 -\log_{10}\!\left[\max\!\left\{
 \frac{c_{(2)}-c_{(1)}}{\max\{\operatorname{IQR}(C_{\mathrm{pred}}),10^{-12}\}},
 10^{-12}\right\}\right].
\end{equation}
The Push-T analysis was exploratory. Leave-one-trajectory-out logistic models
include model identity, with model-specific top-quartile thresholds learned
on the other 19 trajectories. Bootstrap intervals resample saved out-of-fold
predictions without refitting the complete pipeline, and do not account for
choosing this diagnostic after the primary result.

On 562 matched records, unfitted normalized-ambiguity-score AUROC is $.714$ ($[.665,.762]$), compared
with $.706$ for the fitted probability. The flagged group covers 25.3\% of
records and switches 82.4\% of the time versus 51.0\% outside it. However,
flagged-minus-unflagged normalized regret of $a_{\mathrm{pred}}$ is $-.007$ for distance utility
($[-.039,.026]$) and $-.056$ for reward ($[-.143,.026]$), neither establishing
greater physical regret.

Before the secondary Wall extension, model thresholds were frozen at
1.2899081787 for JEPA-WM and 1.4614315032 for DINO-WM, without Wall retuning.
Across all 179 valid records, the normalized score has AUROC $.727$ versus
$.735$ for raw margin. These descriptive estimates provide no normalization
advantage. The signed native gap is
$\min_iD(s_i,s_g)-D(s_{i_{\mathrm{pred}}},s_g)$; larger values mean less distance regret.
Its flagged-minus-unflagged effect is $+.393$
($[-1.704,2.627]$), using the 19 complete cases containing both groups.
Thus prospective freezing limits retuning but does not establish physical-risk
calibration, verification benefit, or an improvement over the simpler baseline.

\section{Artifacts, recomputation, and scope}
\label{app:repro}

The new analysis reads retained candidate scores, actions, terminal states
where available, native distances, and task outcomes. It verifies candidate
identity across models and reproduces the original effects, intervals, tests,
and families before computing the state-selector extension. Its numerical
review independently matched all 36 pooled/model-specific summaries and
intervals, valid-record selector outcomes, and cluster arrays. Original source
artifacts and result files are not overwritten.
These revision materials are retained locally for author release; this work
does not claim that the new package has already been publicly uploaded.

The camera-ready artifact directory
\path{output/task_native_selector_20261002/} contains the checked specification
and implementation reviews, \texttt{summary.csv}, \texttt{row\_effects.csv},
\texttt{cluster\_effects.csv}, \texttt{selection\_disagreements.csv},
\texttt{original\_results.json}, \texttt{physical\_support.json},
\texttt{RESULTS\_INTERPRETATION.md}, reproduction instructions, tests, and a
portable reproduction bundle. \texttt{SCIENTIFIC\_FREEZE.md} records the
pre-manuscript numerical freeze. The numerical summary SHA-256 is
\texttt{3eb5cb00fd7be3f9e01a6f1e6277b003}\allowbreak\texttt{02d0e1549c7d8c8845fba59570bb09c6}.
The reviewed result SHA-256 is
\texttt{f679e348860c2e4db3af4eb98c506b29}\allowbreak\texttt{532127482c2fd60135d8d142801ca8d9}.
Repository and portable relative commands are provided in the respective
\texttt{README.md} files. The portable package was checked from a relocated
directory with Python's isolated mode and an existing NumPy installation;
all declared numerical outputs matched exactly. This is not a fresh dependency
installation or simulator rerun. Python and NumPy suffice for the new CPU analysis;
it does not run models or environments.

Stored native distances are authoritative for retrospective selection.
The compact Push-T export does not contain full initial-state and raw-image
evidence, and exported Wall terminal states have rounded float32 precision;
recomputing native distances from those rounded arrays would change the
numerical object being audited. The retained distances and reward arrays
are sufficient for the reported retrospective comparisons, not for complete
regeneration of the original simulator runs.

The earlier compact Wall bundle contains 179 valid model records with 300
candidates each: 53,700 model--candidate records. Repeated physical sets
across models remain separate prediction records, not independent physical
trials. Its 180-record manifest and failure ledger preserve the failed row.
The Push-T manifest, preregistration, recovery amendment, issuance and attempt
ledgers, analysis outputs, and source bindings are retained separately. The
artifacts support analysis recomputation, not reproduction of model training
or a claim that every raw observation is included in one compact export.

\paragraph{Tool assistance and responsibility.}
LLM agents assisted implementation, analysis scripting, verification, and
drafting. Numerical claims were checked against retained artifacts and an
independent recomputation. The authors remain responsible for study design,
scientific claims, and presentation; agent assessments are not empirical
evidence. These simulator audits provide no deployment or robot-safety guarantee.
